\chapter{IMPLEMENTASI DAN PENGUJIAN SOFTWARE}
\section{Desain}
\subsection{Database}

\vspace{0.5cm}
\noindent\fbox{\begin{minipage}{\textwidth}
\textit{\textbf{Pedoman Penulisan:} Pembuatan desain basis data disesuaikan dengan aplikasi yang akan dibangun oleh penulis.}
\end{minipage}}
\vspace{0.5cm}

\subsection{Software Architecture}

\vspace{0.5cm}
\noindent\fbox{\begin{minipage}{\textwidth}
\textit{\textbf{Pedoman Penulisan:} Pada sub bab ini menjelaskan pseudocode dari algorithma yang digunakan untuk menyelesaikan permasalahan dalam pembangunan aplikasi.}
\end{minipage}}
\vspace{0.5cm}

\subsection{User Interface}

\vspace{0.5cm}
\noindent\fbox{\begin{minipage}{\textwidth}
\textit{\textbf{Pedoman Penulisan:} Pada uraian ini dijelaskan model dari antar muka pengguna (user interface) dari aplikasi yang dikembangkan. Menampilkan desain antar muka form-form yang dibuat.}
\end{minipage}}
\vspace{0.5cm}


\section{Code Generation}

\vspace{0.5cm}
\noindent\fbox{\begin{minipage}{\textwidth}
\textit{\textbf{Pedoman Penulisan:} a. Jika program yang dibuat adalah pemrograman berorientasi objek (OOP), tampilkan hasil generate class-class yang dibuat.\\
b. Jika program yang dibuat adalah pemrograman terstruktur, tampilkan listing program pada form yang terkait proses bisnis utamanya.}
\end{minipage}}
\vspace{0.5cm}


\section{Testing}

\vspace{0.5cm}
\noindent\fbox{\begin{minipage}{\textwidth}
\textit{\textbf{Pedoman Penulisan:} Pengujian (testing) unit yang dilakukan pada unit program yang telah dibuat, boleh menggunakan metode pengujian yang dikuasai (seperti Blackbox testing).}
\end{minipage}}
\vspace{0.5cm}


\section{Support}
\subsection{Publikasi Software}
\subsection{Spesifikasi Hardware dan Software}

\vspace{0.5cm}
\noindent\fbox{\begin{minipage}{\textwidth}
\textit{\textbf{Pedoman Penulisan:} Berisikan dimana software tersebut akan dipublish (web, desktop, android, dll) serta spesifikasi hardware dan software yang akan digunakan untuk pemasangan aplikasi.}
\end{minipage}}
\vspace{0.5cm}

